\section*{Chapter \ref{ch:g10:supply} --- Supply}
% figdata: figdata/grade-10/05-supply/out/ (crude_us, soy_area); Brenland and Kestria numbers invented.

\begin{solution}{exo:g10:supply:1}
Yes: the quantity offered rises with the price. From 1 to 1.4 dalens it rises by $1700 - 1000 = 700$ litres over four tenths of a dalen,
175 litres per tenth on average (200 per tenth between 1 and 1.2, 150 per tenth between 1.2 and 1.4).
\end{solution}

\begin{solution}{exo:g10:supply:2}
(a) Left. (b) No shift: a move along the curve. (c) Right. (d) Left. (e) Right: more trees bearing fruit, that is more capacity.
\end{solution}

\begin{solution}{exo:g10:supply:3}
``Supply rose'': the whole curve moved right, so more is offered at every price (for taxis: the town issues more licences, or a cheaper car
model lowers the cost of a ride). ``The quantity supplied rose'': sellers moved along the same curve because the price rose (for taxis: fares
rise on New Year's Eve and more drivers work that night).
\end{solution}

\begin{solution}{exo:g10:supply:4}
Lowest in 2008, at about 5.0 million barrels a day. Production first passed its 1985 level (8.97 million) in 2015 (9.45 million).
\end{solution}

\begin{solution}{exo:g10:supply:5}
At $p = 3$: only the first seller produces, $3 - 2 = 1$. At $p = 6$: $(6 - 2) + (12 - 8) = 4 + 4 = 8$.
\end{solution}

\begin{solution}{exo:g10:supply:6}
By \cref{prop:g10:supply:shift}, with $q_B = 3(p - 5)$ and $t = 2$: $q_B = 3(p - 7) = 3p - 21$ for $p \ge 7$. At 8 dalens: $24 - 15 = 9$ before,
$24 - 21 = 3$ after (hundreds of loaves): a fall of $3 \times 2 = 6$.
\end{solution}

\begin{solution}{exo:g10:supply:7}
$q_A = 2(p - 2 + 1.5) = 2p - 1$ for $p \ge 0.5$. For $p \ge 5$: $Q = (2p - 1) + (3p - 15) = 5p - 16$. At 9 dalens: $45 - 16 = 29$ against $45 - 19 = 26$
before: 3 hundred loaves more ($2 \times 1.5$, all from A).
\end{solution}

\begin{solution}{exo:g10:supply:8}
After a week: almost no change, since boats and crews are fixed; perhaps a few more trips. After a year: more fishing days, some idle boats back at
sea, some new crews: a moderate rise. After ten years: new boats built and new firms entering, a much larger rise, unless the fish stock cannot
sustain it. The long-run response could be stopped by a fixed limit on catches, a fixed number of licences, or a stock of fish that falls as more
is caught.
\end{solution}

\begin{solution}{exo:g10:supply:9}
Brazil: $45.91 - 11.49 = 34.42$ million hectares, $45.91/11.49 \approx 4.0$ times the 1990 area. United States: $34.82 - 22.87 = 11.95$ million
hectares, $34.82/22.87 \approx 1.52$ times. The response was much larger in Brazil, in proportion as well as in hectares.
\end{solution}

\begin{solution}{exo:g10:supply:10}
For $1 \le p < 3$: $Q = 2p - 2$. For $3 \le p < 6$: $Q = (2p - 2) + (p - 3) = 3p - 5$. For $p \ge 6$: $Q = 3p - 5 + 4p - 24 = 7p - 29$. Kinks at $p = 3$
($Q = 4$) and $p = 6$ ($Q = 13$); below $p = 1$, nothing is offered.
\end{solution}

\begin{solution}{exo:g10:supply:11}
In Brenland, timber, an input of furniture, becomes dearer: every piece of furniture costs more to make, and the supply of furniture shifts left. In
Kestria, timber that can no longer be sold to Brenland must be sold at home, so timber becomes cheaper there: the supply of furniture in Kestria
shifts right.
\end{solution}

\begin{solution}{exo:g10:supply:12}
The model predicts a small first-year response and a larger one later. It would be contradicted if, with the price half as high again for ten
years and nothing else changing, the area and output of the crop rose in the first year and then stayed flat, or did not rise at all although land
and farmers were available and nothing (a rule, a disease, a lack of water) stopped them. Before concluding, check costs: a rise in fertiliser or
labour costs at the same time would shift supply left and hide the response.
\end{solution}

\begin{solution}{exo:g10:supply:13}
Once a well is drilled, most of its cost has been paid; pumping one more barrel costs little, so wells already drilled keep pumping even at a much
lower price, and production falls only a little in the short run (from 9.45 to 8.85 million barrels a day). A new well costs the full price of
drilling and cracking the rock; at the lower price many new wells no longer paid, and drilling slowed. Since the new wells of this kind produce most
of their oil in their first years, the fall in drilling showed up in production only gradually.
\end{solution}

\begin{solution}{exo:g10:supply:14}
Short run: the sum of 200 rising individual curves, each stopping at 60 rides: an upward-sloping curve that becomes vertical at $200 \times 60 =
\num{12000}$ rides a day. Long run: no taxi may enter, so the \omterm{def:g10:supply:short-run}{long-run supply} is the same as the short-run one, with the same vertical part at
\num{12000}. With 100 new licences, both curves become the sum of 300 curves: at each price the quantity is one and a half times as large, and the
vertical part moves to \num{18000} rides a day.
\end{solution}

\begin{solution}{exo:g10:supply:15}
The price rises by $10/40 = 25\,\%$. Short run: the quantity rises by 10\,\%, so $10/25 = 0.4\,\%$ for each 1\,\% of price. Long run: 40\,\%, so
$40/25 = 1.6\,\%$ for each 1\,\%. The long-run number is larger because, given time, firms add capacity and new firms enter, so the higher price calls
forth units that could not be produced at all with the old capacity.
\end{solution}

\begin{solution}{pb:g10:supply:1}
Quantities are in thousands of litres a year; prices in kests per hundred litres.

\textbf{1.} Small farms 20 kests, large farms 15 kests.

\textbf{2.} Small: $2 \times (30 - 20) = 20$. Large: $10 \times (30 - 15) = 150$.

\textbf{3.} For $15 \le p < 20$ only the 10 large farms produce: $Q = 10 \times 10(p - 15) = 100p - 1500$. For $p \ge 20$: $Q = 100p - 1500 + 40 \times
2(p - 20) = 100p - 1500 + 80p - 1600 = 180p - 3100$.

\textbf{4.} At $p = 20$, where the small farms start: $Q = 2000 - 1500 = 500$. At 30 kests: $Q = 5400 - 3100 = 2300$.

\textbf{5.} The large farms: $10 \times 150 = 1500$ of the 2300 (about 65\,\%), the small ones $40 \times 20 = 800$. The large farms start at a lower
price and each one's quantity rises five times as fast with the price.

\textbf{6.} Small: $q = 2(p - 25)$ for $p \ge 25$. Large: $q = 10(p - 20)$ for $p \ge 20$ (\cref{prop:g10:supply:shift}).

\textbf{7.} For $p \ge 25$: $Q = 80(p - 25) + 100(p - 20) = 180p - 4000$. At 30 kests: $5400 - 4000 = 1400$.

\textbf{8.} $2300 - 1400 = 900$, a fall of about 39\,\%. It is a shift: the world price did not change; the cost of every unit did.

\textbf{9.} $180p - 4000 = 2300$ gives $p = 6300/180 = 35$ kests: the curve has moved up by exactly the 5 kests of the cost rise.

\textbf{10.} $Q = 180 \times 36 - 3100 = 6480 - 3100 = 3380$.

\textbf{11.} $3380 - 2300 = 1080$, that is $1080/2300 \approx 47.0\,\%$.

\textbf{12.} Small farms: $40 \times 2 \times 6 = 480$; large farms: $10 \times 10 \times 6 = 600$. Total 1080.

\textbf{13.} A move along the \omterm{def:g10:supply:short-run}{short-run supply} curve: the same 50 farms respond to a higher price.

\textbf{14.} $10 \times (36 - 15) = 210$.

\textbf{15.} $Q = 180p - 3100 + 15 \times 10(p - 15) = 180p - 3100 + 150p - 2250 = 330p - 5350$.

\textbf{16.} $330 \times 36 - 5350 = 11880 - 5350 = 6530$.

\textbf{17.} $6530 - 2300 = 4230$, almost three times today's production ($6530/2300 \approx 2.84$).

\textbf{18.} $4230/1080 \approx 3.9$ times.

\textbf{19.} $4230/6 = 705$ thousand litres per kest, against 180 for the short-run \omterm{def:g10:supply:market-supply}{market supply}: the long-run response is about four times as steep in
quantity (the long-run curve, drawn with the price on the vertical axis, is four times flatter).

\textbf{20.} After the world price rose from 30 to 36 kests and stayed there, Kestria's milk production rose by \num{1080} thousand litres a year in the
first year and by \num{4230} thousand litres once new farms had opened, about 3.9 times the first-year rise.
\end{solution}
